% Checking notes accompanying the paper; equation references are to the paper.
\section{The formal character bridge}
This section gives the definitions and direct formal proof of the character formula.
The accompanying Lean project defines $S_n$ as the permutations of
$\{0,\ldots,n-1\}$ and works in its integral monoid algebra.
The class sum includes each permutation of the prescribed cycle type
with coefficient one. The matrix entry is the coefficient of a chosen
permutation of type $\mu$ in $K_{(n)}K_\nu$; the development proves
that the choice has the required cycle type and that these entries
give the expansion~\eqref{eq:matrix}.

In ASCII Lean notation, the final theorem is:
\begin{verbatim}
theorem stanley_class_sum_snf (n : Nat) (hn : 0 < n) :
    And (HasStanleySNF n (M n))
      (forall r, r + 1 < n ->
        Dvd.dvd (stanleyD n r) (stanleyD n (r + 1)))
\end{verbatim}
Here \texttt{stanleyD} is~\eqref{eq:d}, with its division proved exact.
The predicate \texttt{HasStanleySNF} asserts the existence of a finite
index set $\kappa$ and an integral basis indexed by
$\{0,\ldots,n-1\}\sqcup\kappa$, for which the image is the span of
$d_r b_r$. This is precisely~\eqref{eq:image}. Positivity and
independence turn that span into the direct sum stated in the theorem.
The written basis-change argument in Section~\ref{sec:assembly} identifies it with the
usual diagonal Smith form. In particular the formal statement does not
replace the group-ring matrix by a matrix defined through a character
formula.

For the character formula, the formal proof uses an explicit model
of the exterior algebra of $\Q^n$, with basis indexed by subsets.
Permutations act by signed permutations of these basis vectors.
Exterior multiplication by $e=e_1+\cdots+e_n$ and contraction by
the all-ones linear functional give operators $E$ and $I$ satisfying
$IE+EI=n$. Thus $P=IE/n$ is an equivariant projection onto
$\ker I$. Taking traces of the degree-$k$ action on this image gives
the coefficients in~\eqref{eq:hook}.

In the conjugacy average below, $a,b$ are arbitrary permutations in $S_n$,
and $\chi_k$ is evaluated on their cycle types.
For the interior degrees $1\le k<n$, the degree-$k$ commutant is
determined by the signed action on pairs
of $k$-subsets: a transposition in their difference kills an entry when
their intersection has size at most $k-2$, and the remaining orbits
leave a two-dimensional span of the identity and the projection $P$.
Here the matrix entries are indexed by pairs of subsets: a transposition
of two elements belonging to one subset but not the other fixes one
subset and reverses the sign on the other.
Restriction to the image of $P$ makes every commuting operator scalar.
At degree $0$, the representation is one-dimensional and $P$ is the
identity, so the same scalar conclusion holds directly.
Taking traces of conjugacy averages consequently gives
\[
 \sum_{g\in S_n}\chi_k(agb g^{-1})
       =\frac{n!}{\binom Nk}\chi_k(a)\chi_k(b).
\]
Evaluation of~\eqref{eq:hook} at $q=-1$ gives the indicator identity
$\sum_k(-1)^k\chi_k(\mu)=n\,\mathbf1_{\mu=(n)}$.
This identity, the preceding conjugacy average, and orbit--stabilizer
prove Proposition~\ref{prop:bridge} directly in Lean. Consequently,
the formal proof does not take the classical representation-theoretic
theorems used in its shorter written derivation as extra hypotheses.

The remaining development follows the polynomial pairing, the two
lattice lemmas, row divisibility, and assembly above. It uses Lean's
kernel and the mathlib library; no certificate generator, numerical
approximation, or external computational oracle is required. From the
bundle's \texttt{lean/} directory, the standard commands are
\begin{verbatim}
lake exe cache get
lake build
\end{verbatim}
The build includes checks of the final statement and its logical
dependencies. The accompanying Lean guide documents these checks and
the correspondence of definitions. The paper can be rebuilt from its
source with \texttt{pdflatex}.

\end{appendices}


% Exact final-theorem correspondence and permitted logical dependencies.
\section{Formal verification}\label{sec:formal}
We state exactly the formal endpoint and its mathematical interpretation.
In the namespace \texttt{StanleySnf}, the theorem
\texttt{stanley\_class\_sum\_\allowbreak snf} states the following. For every natural number $n$,
if $0<n$, then \texttt{HasStanleySNF n (M n)} and, for every natural
number $r$, $r+1<n$ implies
\texttt{stanleyD n r} divides \texttt{stanleyD n (r+1)}.
There are no other hypotheses.
The predicate \texttt{HasStanleySNF n A} means that there exist a type
$\kappa$, a finite enumeration of $\kappa$, and a $\Z$-basis
$b$ of the integer-valued functions on partitions of $n$, indexed by
$\texttt{Fin }n\sqcup\kappa$, such that the range of multiplication of
column vectors by $A$ equals the $\Z$-span of
$\{\texttt{stanleyD n r}\,b(\mathrm{inl}\,r):r\in\texttt{Fin }n\}$.
The natural-number function \texttt{stanleyD} is exactly~\eqref{eq:d};
its division is proved exact and its relevant values positive.
Thus the span is the direct sum in~\eqref{eq:image}.

The permutations are those of $\texttt{Fin }n=\{0,\ldots,n-1\}$,
including fixed points in their partition; $n=0$ is excluded.
The entry \texttt{M n} is defined in the integral monoid algebra by the
coefficient at a chosen permutation of type $\mu$ in $K_{(n)}K_\nu$.
The fidelity theorems \texttt{Fid.classSum\_coeff},
\texttt{Fid.rep\_spec}, \texttt{Fid.parts\_eq\_iff\_isConj}, and
\texttt{Fid.M\_structure} respectively assert that a class-sum coefficient
is one exactly at the specified partition and zero otherwise, that the
representative has that partition, that partition equality is equivalent
to conjugacy to the representative, and, for $0<n$, that
$K_{(n)}K_\nu=\sum_\mu M_n[\mu,\nu]K_\mu$ for every $\nu\vdash n$.
These fix the row/output and column/input conventions.

The theorem \texttt{bridge} says: for every natural $n$ with $0<n$ and
every $\mu,\nu\vdash n$, the rational cast of \texttt{M n \ensuremath{\mu} \ensuremath{\nu}}
equals \texttt{bridgeFormula n \ensuremath{\mu} \ensuremath{\nu}}.
That formula is precisely the hook sum above.
The intermediate theorem \texttt{main} has the same conclusion as
\texttt{stanley\_class\_sum\_\allowbreak snf}, with the additional hypothesis that
this bridge equality holds for every $\mu,\nu$; \texttt{bridge} discharges it.

The permitted axioms are exactly \texttt{propext},
\texttt{Classical.choice}, and \texttt{Quot.sound}; the default build guards
their reports. Trust rests on Lean's kernel and the pinned Lean/mathlib
inputs, not an external oracle. The conversion of the image-lattice
statement to the usual diagonal form, the count $p(n)-n$ of zero entries,
and Corollary~\ref{cor:primes} are written deductions, not separate Lean
endpoints. The main text gives their proofs.
From \texttt{lean/}, reproduce the formal check with
\texttt{lake exe cache get} followed by \texttt{lake build};
\texttt{lean/README.md} gives the pins and checks.

